% =============================================================================
% INCOHERENTIA: LA TRANSMISIÓN (SPANISH ONLY)
% Status: DRAFTED
% Placement: Between XX and XXI
% Fully Spanish narrative. Aisha and Saladin at the Italian port.
% Imperfect copying as survival strategy. No translation provided.
% Core theme: arrival as transformation, not preservation.
% =============================================================================

\chapter*{La Transmisión}
\addcontentsline{toc}{chapter}{La Transmisión}

\setstretch{1.03}

No fue el texto lo que viajó.

O no solamente.

Lo que se desplazaba de un lugar a otro no era una obra terminada, ni un
argumento cerrado, sino algo más inestable: una forma de continuar hablando
sin garantizar el mismo sentido.

En el puerto, las cajas se abrían con una mezcla de prisa y descuido. Los
nombres se repetían en voz alta, pero cada repetición los alteraba ligeramente.

---¿Esto dice ``Ibn Rushd''? ---preguntó uno de los cargadores.

---Eso parece ---respondió otro---. O algo parecido.

El escriba tomó el documento, frunció el ceño, y volvió a trazar las letras
con tinta fresca. No corregía. Ajustaba.

A unos pasos de distancia, Aisha observaba sin intervenir. Había aprendido
que la transmisión no dependía de la fidelidad, sino de la continuidad. Un
error no detenía el proceso. A veces lo hacía posible.

Saladin estaba junto a ella, mirando las cajas como si intentara entender qué
parte de aquello realmente importaba.

---Si cambian los nombres ---dijo---, ¿no cambia también lo que significan?

Aisha tardó en responder.

---Sí ---dijo finalmente---. Pero no necesariamente deja de funcionar.

Un hombre pasó junto a ellos cargando un volumen mal cerrado. Algunas hojas
se deslizaron y cayeron al suelo. Nadie se detuvo de inmediato. El hombre
siguió caminando. Otro recogió una de las hojas, la miró brevemente, y la
colocó en la pila equivocada.

Aisha lo vio. No dijo nada.

---¿No vas a corregirlo? ---preguntó Saladin.

Ella negó.

---No puedo corregir todo ---dijo---. Y no todo necesita ser corregido.

---Entonces, ¿cómo sabes qué importa?

Aisha lo miró por primera vez directamente.

---No lo sé ---dijo---. Lo descubro viendo qué sobrevive.

El ruido del puerto aumentó a medida que el día avanzaba. Voces en diferentes
lenguas se cruzaban sin necesidad de coincidir. Nada completamente estable.
Nada completamente perdido.

Saladin se inclinó y recogió una hoja que el viento había empujado hacia sus
pies. La sostuvo con cuidado.

---¿Puedes leer esto? ---preguntó.

Aisha se acercó. Observó la página. Las palabras le resultaban familiares,
pero no exactas. Habían cambiado. O quizá ella había cambiado.

---Sí ---dijo---. Creo que sí.

---¿Y significa lo mismo?

Aisha dudó. Luego negó suavemente.

---No exactamente.

Saladin asintió, como si esa fuera la única respuesta que tenía sentido.

Doblando la hoja, la colocó de nuevo entre las otras. No en su lugar original.
En uno posible.

El trabajo continuó. Las cajas se cerraron. Los nombres fueron escritos otra
vez. Y otra vez cambiaron.

Al final del día, nada estaba exactamente donde había comenzado.

Y sin embargo, nada se había detenido.

Eso, pensó Aisha, era suficiente.
